\documentclass[border=10pt]{standalone}
\usepackage{tikz,ctex}
\usetikzlibrary{positioning}

\begin{document}

\begin{tikzpicture}[
    % 全局设置
    node distance = 1.6cm and 1.4cm,
    every node/.style = {
        draw,
        rounded corners,
        thick,
        align = center,
        font = \sffamily\small,
        minimum width = 2.0cm,
        minimum height = 0.9cm
    },
    % 样式定义
    centerStyle/.style = {
        fill = blue!45,
        text = white,
        font = \sffamily\bfseries\large,
        line width = 2pt,
        minimum size = 2.8cm,
        circle
    },
    branchTitle/.style = {
        font = \sffamily\bfseries,
        line width = 1.5pt
    },
    subNode/.style = {
        fill = white,
        font = \sffamily\small,
        minimum width = 2.2cm,
        minimum height = 0.8cm
    },
    leafNode/.style = {
        fill = white,
        font = \sffamily\footnotesize,
        minimum width = 1.9cm,
        minimum height = 0.7cm
    },
    arrowStyle/.style = {
        ->,
        thick,
        >=stealth,
        shorten >=2pt
    }
]

    % ================= 1. 中心节点 =================
    \node[centerStyle] (center) {9.4.参数共享 \\ Parameter Sharing};

    % ================= 2. 上方分支：硬权重共享 (红) =================
    \node[branchTitle, fill=red!15, above=of center] (hard) {硬权重共享 \\ Hard Sharing};

    \node[subNode, above=of hard, xshift=-6cm] (tie) {参数绑定 \\ Parameter Tying};
    \node[subNode, above=of hard, xshift=-2cm] (dof) {自由度降低 \\ Fewer DOF};
    \node[subNode, above=of hard, xshift=2cm] (bias) {归纳偏置 \\ Inductive Bias};
    \node[subNode, above=of hard, xshift=6cm] (cnn) {卷积网络 \\ CNN};

    % ================= 3. 右侧分支：软权重共享 (蓝) =================
    \node[branchTitle, fill=orange!15, right=of center] (soft) {软权重共享 \\ Soft Sharing};

    \node[subNode, right=of soft, yshift=4cm] (mix) {高斯混合 \\ Gaussian Mixture};
    \node[subNode, right=of soft, yshift=2cm] (mean) {均值 $\mu_j$ \\ Means};
    \node[subNode, right=of soft, yshift=0cm] (var) {方差 $\sigma_j^2$ \\ Variances};
    \node[subNode, right=of soft, yshift=-2cm] (mixcoef) {混合系数 $\pi_j$ \\ Mixing Coefficients};
    \node[subNode, right=of soft, yshift=-4cm] (post) {后验概率 \\ Posterior};

    % ================= 4. 下方分支：学习与优化 (绿) =================
    \node[branchTitle, fill=green!15, below=of center] (learn) {学习与优化 \\ Learning \& Optimization};

    \node[subNode, below=of learn, xshift=-6cm] (grad) {梯度下降 \\ Gradient Descent};
    \node[subNode, below=of learn, xshift=-2cm] (softmax) {Softmax 参数化 \\ Softmax};
    \node[subNode, below=of learn, xshift=2cm] (exp) {对数方差 \\ Log-Variance};
    \node[subNode, below=of learn, xshift=6cm] (joint) {联合优化 \\ Joint Optimization};

    % ================= 5. 左侧分支：扩展应用 (紫) =================
    \node[branchTitle, fill=purple!15, left=of center] (ext) {扩展应用 \\ Extensions};

    \node[subNode, left=of ext, yshift=3cm] (gen) {生成模型 \\ Generative Model};
    \node[subNode, left=of ext, yshift=1cm] (disc) {判别模型 \\ Discriminative Model};
    \node[subNode, left=of ext, yshift=-1cm] (hybrid) {混合方法 \\ Hybrid Approach};
    \node[subNode, left=of ext, yshift=-3cm] (unlabel) {无标签数据 \\ Unlabelled Data};

    % ================= 连线逻辑 =================
    % 中心 -> 四大分支标题
    \draw[arrowStyle, red] (center) -- (hard);
    \draw[arrowStyle, blue] (center) -- (soft);
    \draw[arrowStyle, green!60!black] (center) -- (learn);
    \draw[arrowStyle, purple] (center) -- (ext);

    % 分支标题 -> 子节点 (上方)
    \draw[arrowStyle, red!60] (hard) -- (tie.south);
    \draw[arrowStyle, red!60] (hard) -- (dof.south);
    \draw[arrowStyle, red!60] (hard) -- (bias.south);
    \draw[arrowStyle, red!60] (hard) -- (cnn.south);

    % 分支标题 -> 子节点 (右侧)
    \draw[arrowStyle, blue!60] (soft) -- (mix.west);
    \draw[arrowStyle, blue!60] (soft) -- (mean.west);
    \draw[arrowStyle, blue!60] (soft) -- (var.west);
    \draw[arrowStyle, blue!60] (soft) -- (mixcoef.west);
    \draw[arrowStyle, blue!60] (soft) -- (post.west);

    % 分支标题 -> 子节点 (下方)
    \draw[arrowStyle, green!60!black] (learn) -- (grad.north);
    \draw[arrowStyle, green!60!black] (learn) -- (softmax.north);
    \draw[arrowStyle, green!60!black] (learn) -- (exp.north);
    \draw[arrowStyle, green!60!black] (learn) -- (joint.north);

    % 分支标题 -> 子节点 (左侧)
    \draw[arrowStyle, purple!60] (ext) -- (gen.east);
    \draw[arrowStyle, purple!60] (ext) -- (disc.east);
    \draw[arrowStyle, purple!60] (ext) -- (hybrid.east);
    \draw[arrowStyle, purple!60] (ext) -- (unlabel.east);

\end{tikzpicture}

\end{document}